\documentclass[12pt,a4paper]{article}
\usepackage[utf8]{inputenc}
\usepackage[polish]{babel}
\usepackage[T1]{fontenc}
\usepackage[margin=20mm,top=10mm,bottom=20mm]{geometry}
\usepackage{setspace}
\usepackage{graphicx}
\usepackage{array}
\usepackage{booktabs}
\usepackage{multirow}
\usepackage{listings}
\usepackage{xcolor}
\usepackage{hyperref}
\usepackage{fancyhdr}
\usepackage{indentfirst}
\usepackage{amsmath}

% Ustawienie odstępu między wierszami 1.5
\onehalfspacing

% Ustawienie czcionki Times New Roman (wymagane przez UZ)
\usepackage{mathptmx}

% Nagłówek i stopka
\pagestyle{fancy}
\fancyhf{}
\rhead{\thepage}
\lhead{Projekt BREW - Dokumentacja}
\renewcommand{\headrulewidth}{0.4pt}

% Konfiguracja listingów (kod C)
\lstset{
    language=C,
    basicstyle=\ttfamily\small,
    keywordstyle=\color{blue},
    commentstyle=\color{gray},
    stringstyle=\color{red},
    breaklines=true,
    showstringspaces=false,
    frame=single,
    backgroundcolor=\color{gray!10}
}

% Wcięcie akapitu 0.6cm (4 pozycje)
\setlength{\parindent}{0.6cm}

\title{\textbf{INTELIGENTNY SYSTEM KONTROLI I MONITOROWANIA\\PROCESU WARZENIA PIWA DOMOWEGO}}
\author{Jacek Kuczak \and Kornel Niewiadomski}
\date{23 stycznia 2026 r.}

\begin{document}

% --- STRONA TYTUŁOWA ---
\thispagestyle{empty}
\begin{center}
    {\Large \textbf{UNIWERSYTET ZIELONOGÓRSKI}} \\
    \vspace{0.5cm}
    {\Large \textbf{WYDZIAŁ INFORMATYKI, ELEKTROTECHNIKI I AUTOMATYKI}} \\
    \vspace{3cm}
    {\huge \textbf{DOKUMENTACJA PROJEKTOWA}} \\
    \vspace{1cm}
    {\Large \textbf{INTELIGENTNY SYSTEM KONTROLI I MONITOROWANIA\\PROCESU WARZENIA PIWA DOMOWEGO}} \\
    \vspace{0.5cm}
    {\large (BREW - Beer Brewing Monitoring System)} \\
    \vspace{3cm}
    \begin{minipage}{0.4\textwidth}
        \begin{flushleft} \large
            \textbf{Autorzy:}\\
            Jacek Kuczak\\
            Kornel Niewiadomski
        \end{flushleft}
    \end{minipage}
    \begin{minipage}{0.4\textwidth}
        \begin{flushright} \large
            \textbf{Prowadzący:}\\
            [Imię Nazwisko Prowadzącego]
        \end{flushright}
    \end{minipage}
    \vfill
    {\large Zielona Góra, 2026}
\end{center}
\newpage

% --- SPIS TREŚCI ---
\tableofcontents
\newpage

% ============================================================================
\section{Wprowadzenie}
% ============================================================================
Projekt BREW stanowi inteligentny system automatyzacji i monitorowania procesu warzenia oraz fermentacji piwa domowego. System wykorzystuje trzy mikrokontrolery ESP32-C3 Mini komunikujące się za pomocą Bluetooth Low Energy (BLE) oraz Wi-Fi. Głównym celem jest zapewnienie piwowarowi domowemu precyzyjnych danych o temperaturze, gęstości oraz aktywności fermentacji w czasie rzeczywistym.

% ============================================================================
\section{Tytuł projektu i cel}
% ============================================================================
\textbf{Pełna nazwa:} Inteligentny system kontroli i monitorowania procesu warzenia piwa domowego za pomocą mikrokontrolerów ESP32 z wykorzystaniem BLE i czujników pomiarowych.

\textbf{Cel projektu:} Zbudowanie modularnego systemu monitorowania, który umożliwi automatyczne dozowanie wody, pomiar gęstości bez otwierania fermentora oraz śledzenie postępów fermentacji na urządzeniach mobilnych.

% ============================================================================
\section{Specyfikacja S.M.A.R.T.}
% ============================================================================
\begin{itemize}
    \item \textbf{Specific (Konkretny):} Stworzenie trzech niezależnych modułów (Przelewowy, Sensoryczny, Monitorujący) zintegrowanych w jedną sieć BLE.
    \item \textbf{Measurable (Mierzalny):} Uzyskanie dokładności pomiaru przepływu wody $\pm 3\%$, temperatury $\pm 0.5^\circ C$ oraz gęstości $\pm 0.001$ SG.
    \item \textbf{Achievable (Osiągalny):} Wykorzystanie popularnych komponentów (ESP32, DS18B20, MPU6050) i sprawdzonych bibliotek programistycznych.
    \item \textbf{Relevant (Istotny):} System rozwiązuje realny problem piwowarów – konieczność otwierania fermentora do pomiaru gęstości, co grozi infekcją piwa.
    \item \textbf{Time-bound (Terminowy):} Realizacja projektu w ciągu 12-13 tygodni, z terminem oddania 23 stycznia 2026 r.
\end{itemize}

% ============================================================================
\section{Podział obowiązków}
% ============================================================================
\subsection{Jacek Kuczak -- System przelewowy}
Odpowiada za warstwę fizyczną dozowania wody i interfejs użytkownika na urządzeniu stacjonarnym.
\begin{itemize}
    \item Implementacja sterowania elektrozaworem (BREW-1).
    \item Programowanie algorytmu przepływomierza (BREW-2).
    \item Obsługa klawiatury UART (BREW-3) i wyświetlacza OLED (BREW-5).
\end{itemize}

\subsection{Kornel Niewiadomski -- Sensoryka i Komunikacja}
Odpowiada za zbieranie danych z czujników zanurzeniowych oraz przesyłanie ich bezprzewodowo.
\begin{itemize}
    \item Obsługa czujników temperatury i akcelerometru (BREW-6, BREW-7).
    \item Implementacja stosu komunikacyjnego BLE/Wi-Fi (BREW-10).
    \item Kalibracja algorytmu gęstości grawitacyjnej (BREW-9).
\end{itemize}

% ============================================================================
\section{Rzeczywisty harmonogram realizacji}
% ============================================================================
\begin{table}[h!]
\centering
\begin{tabular}{|l|c|c|c|}
\hline
\textbf{Zadanie} & \textbf{Plan (Tyg.)} & \textbf{Rzeczywistość (Tyg.)} & \textbf{Odchylenie} \\ \hline
Zamówienie części & 1-2 & 1-2 & 0 \\ \hline
Prototypowanie (BREW-1,2,3) & 3-5 & 3-5 & 0 \\ \hline
System zasilania (BREW-4) & 5-6 & 5-7 & +1 tydz. \\ \hline
Komunikacja BLE (BREW-10) & 8-10 & 8-11 & +1 tydz. \\ \hline
Testy i Dokumentacja & 11-12 & 11-13 & 0 \\ \hline
\end{tabular}
\caption{Harmonogram realizacji projektu BREW}
\end{table}

% ============================================================================
\section{Opis realizacji zadań -- Jacek Kuczak}
% ============================================================================
\subsection{BREW-2: Obsługa czujnika przepływu}
Czujnik YF-S201 generuje impulsy, których częstotliwość $f$ zależy od natężenia przepływu $Q$.
Wzór na przepływ w litrach na minutę:
\begin{equation}
Q = \frac{f}{7.5}
\end{equation}
Całkowita objętość $V$ jest obliczana jako całka z przepływu po czasie:
\begin{equation}
V = \sum_{t=0}^{n} Q(t) \cdot \Delta t
\end{equation}

\begin{lstlisting}[caption=Kod obsługi przerwania przepływomierza]
void IRAM_ATTR flow_sensor_isr(void* arg) {
    pulse_count++;
}
// Obliczanie co 1 sekunde:
float flow_rate = pulse_count / 7.5;
total_volume += flow_rate / 60.0;
\end{lstlisting}

% ============================================================================
\section{Opis realizacji zadań -- Kornel Niewiadomski}
% ============================================================================
\subsection{BREW-9: Kalibracja gęstości (Tilt)}
Gęstość (SG) obliczana jest na podstawie kąta nachylenia $\theta$ czujnika zanurzonego w piwie.
Zastosowano wielomian 3. stopnia:
\begin{equation}
SG = a\theta^3 + b\theta^2 + c\theta + d
\end{equation}
Współczynniki $a, b, c, d$ zostały wyznaczone podczas kalibracji w roztworach cukrowych o znanym stężeniu.

\begin{lstlisting}[caption=Implementacja BLE GATT Notify]
uint8_t temp_data[4];
memcpy(temp_data, &current_temp, 4);
esp_ble_gatts_send_indicate(gatts_if, conn_id, char_handle, 
                            sizeof(temp_data), temp_data, false);
\end{lstlisting}

% ============================================================================
\section{Dokumentacja testów i screeny}
% ============================================================================
\begin{itemize}
    \item \textbf{Test BREW-10:} Weryfikacja zasięgu BLE w warunkach domowych (przez ściany). Wynik: stabilne połączenie do 15 metrów.
    \item \textbf{Test BREW-2:} Kalibracja przepływomierza przy użyciu naczynia pomiarowego 5L. Wynik: błąd $\pm 2.4\%$.
\end{itemize}
\textit{[Miejsce na zdjęcie: zmontowany układ sterownika przepływu]}
\textit{[Miejsce na zdjęcie: wykres zmian temperatury z aplikacji mobilnej]}

% ============================================================================
\section{Podsumowanie i wnioski}
% ============================================================================
Projekt został zrealizowany zgodnie z założeniami funkcjonalnymi. Największym wyzwaniem okazała się stabilizacja zasilania dla trzech modułów oraz kalibracja akcelerometru pod kątem gęstości płynu. 
\begin{itemize}
    \item System pozwala na realne oszczędności czasu podczas warzenia.
    \item Komunikacja BLE jest energooszczędna i wystarczająca dla potrzeb browaru domowego.
    \item Planowany rozwój: dodanie modułu sterowania grzałką (PID).
\end{itemize}

% ============================================================================
\section{Bibliografia}
% ============================================================================
\begin{thebibliography}{9}
\bibitem{esp32} Espressif Systems, \textit{ESP32-C3 Technical Reference Manual}, 2025.
\bibitem{ispindel} iSpindel Project, \textit{DIY Tilt Hydrometer Documentation}, GitHub 2024.
\bibitem{brewery} Palmer J., \textit{How to Brew}, Brewers Publications, 2017.
\end{thebibliography}

\end{document}
